\section{Introduction}
\label{sec:intro}

A practitioner recommending a game to a group does not ask ``is this a good game?'' ---
they ask ``is this the right game for these players, right now?'' That question requires
not a quality score but a profile: a characterization of what the game emphasizes,
what kind of experience it produces, and how it differs from the alternatives already
known to the group. TIGER is a five-dimensional framework designed to answer that question.

\subsection{The Limits of Quality-Based Frameworks}

Existing board game evaluation frameworks, from BoardGameGeek's community weight and
rating system to academic complexity indices~\cite{woods2012eurogames,
zagal2008collaborative}, are organized around quality or complexity as a single dominant
axis. This organization is appropriate for ranking but inappropriate for recommendation,
design guidance, or market analysis. Two games can share an identical BGG rating while
occupying opposite corners of design space --- one relying on tight resource tension and
another on rich player interaction --- and a quality-only framework cannot distinguish them.

\subsection{TIGER as a Distinctiveness Framework}

TIGER (Tension, Interaction, Game-Architecture, Elegance, Range) emerges from PCA on
the within-game normalized TIGRIS corpus~\cite{dellaLibera2026pca} and is formalized in
this paper with anchor descriptors at 0, 5, and 10 for each dimension. The framework
does not replace quality assessment --- it supplements it, characterizing the design
\emph{space} a game occupies rather than its position on a quality \emph{ladder}.
